\lampiransection{Lembar Observasi GrabFood dan Kompetitor}
\label{app:observasi-lanjutan-grabfood}

Lampiran ini digunakan untuk mencatat observasi terhadap toko Nasi Gerilya dan kompetitor pada GrabFood. Setelah data awal tanggal 13 Juni 2026 serta 19 dan 20 Juni 2026 tersedia, lembar ini memuat ringkasan observasi yang sudah masuk sekaligus ruang observasi lanjutan untuk triangulasi waktu.

\begin{table}[!htbp]
    \centering
    \caption{Rekap Observasi Awal Toko Nasi Gerilya pada GrabFood}
    \label{tab:rekap-observasi-awal-toko-grabfood}
    \footnotesize
    \setlength{\tabcolsep}{4pt}
    \renewcommand{\arraystretch}{0.95}
    \begin{tabularx}{\textwidth}{|Z{3cm}|Y|Y|}
        \toprule
        \textbf{Aspek Observasi} & \textbf{Data Nasi Gerilya} & \textbf{Kode Temuan / Catatan} \\ \hline
        \midrule
        Tanggal observasi & 13 Juni 2026 & Bukti visual pada Lampiran~\ref{app:evidence-observasi-grabfood-20260613}. \\ \hline
        Identitas dan reputasi & Nasi Gerilya - Glugur Kota; rating 4,7 dari 4.584 penilaian; jarak 6,6 km; status saat screenshot tutup dan buka besok. & TO-GF-01 \\ \hline
        Ringkasan ulasan agregat & Rasa enak, porsi besar, bahan segar, kemasan rapi, menu beragam, dan harga terjangkau. & TO-GF-02 \\ \hline
        Promo aktif & 19 promo/voucher terlihat, termasuk diskon persentase, potongan nominal, diskon ongkir, bank, dan \textit{group order}. & TO-GF-03 \\ \hline
        Menu terlaris & Nasi Padang Rendang Sapi, Nasi Padang Ayam Goreng Bumbu, dan Nasi Padang Ayam Rendang. & TO-GF-04 \\ \hline
        Rentang harga & Keseluruhan Rp1.100--Rp275.000; item tersedia Rp1.100--Rp72.000; paket nasi Padang tersedia Rp27.500--Rp72.000. & TO-GF-05 \\ \hline
        Ketersediaan menu & 6 menu atau produk terlihat berstatus \textit{Nggak tersedia}. & TO-GF-06 \\ \hline
        Ulasan positif & Rasa, porsi, kesegaran bahan, kebersihan, pengemasan, request terpenuhi, dan pelanggan rutin. & TO-GF-07 \\ \hline
        Ulasan negatif/netral & Salah item, kelengkapan pesanan, add-on tidak diberikan, persepsi harga/porsi, kualitas makanan, koordinasi staf, dan pengalaman pascakonsumsi. & TO-GF-08 \\ \hline
        \bottomrule
    \end{tabularx}
    \tablesource{Diolah peneliti berdasarkan observasi GrabFood Nasi Gerilya tanggal 13 Juni 2026.}
\end{table}

\begin{table}[!htbp]
    \centering
    \caption{Rekap Perbandingan Awal Kompetitor pada GrabFood}
    \label{tab:rekap-perbandingan-kompetitor-grabfood}
    \scriptsize
    \setlength{\tabcolsep}{2pt}
    \renewcommand{\arraystretch}{0.88}
    \begin{tabularx}{\textwidth}{|Z{2.8cm}|Z{2.5cm}|Z{2.9cm}|Y|}
        \toprule
        \textbf{Merchant} & \textbf{Reputasi dan Jarak} & \textbf{Rentang Harga} & \textbf{Catatan Analisis Awal} \\ \hline
        \midrule
        Nasi Gerilya - Glugur Kota & \begin{tabular}[t]{@{}l@{}}Rating 4,7;\\4.584 penilaian;\\jarak 6,6 km\end{tabular} & \begin{tabular}[t]{@{}l@{}}Item tersedia\\Rp1.100--\\Rp72.000;\\paket nasi\\Rp27.500--\\Rp72.000\end{tabular} & Jumlah penilaian lebih besar daripada Paripurna dan Dendeng Batokok, tetapi masih di bawah Pondok Gurih, Garuda, dan Istana Krakatau dari sisi akumulasi penilaian. \\ \hline
        Restoran Paripurna - Sei Sikambing B & \begin{tabular}[t]{@{}l@{}}Rating 4,7;\\1.264 penilaian;\\jarak 3,0 km\end{tabular} & \begin{tabular}[t]{@{}l@{}}Rp5.000--\\Rp66.000;\\nasi utama\\Rp24.300--\\Rp66.000\end{tabular} & Lebih dekat dari titik observasi dan rentang menu utama mendekati Nasi Gerilya. \\ \hline
        Rumah Makan Dendeng Batokok - Glugur Darat II & \begin{tabular}[t]{@{}l@{}}Rating 4,8;\\1.479 penilaian;\\jarak 7,3 km\end{tabular} & \begin{tabular}[t]{@{}l@{}}Rp5.000--\\Rp36.000;\\nasi/lauk\\Rp7.000--\\Rp36.000\end{tabular} & Rating lebih tinggi dan batas harga atas lebih rendah, sehingga kuat sebagai pembanding harga bawah. \\ \hline
        Rumah Makan Istana Krakatau - Centre Point Mall & \begin{tabular}[t]{@{}l@{}}Rating 4,6;\\6.256 penilaian;\\jarak 8,4 km\end{tabular} & \begin{tabular}[t]{@{}l@{}}Rp9.000--\\Rp250.000;\\nasi/kotak\\Rp22.500--\\Rp83.000\end{tabular} & Rating lebih rendah, tetapi jumlah penilaian dan spektrum menu premium lebih besar sehingga tetap menjadi pembanding reputasi dan harga. \\ \hline
        RM. Pondok Gurih - Brigjend Katamso & \begin{tabular}[t]{@{}l@{}}Rating 4,9;\\56.157 penilaian;\\jarak 5,0 km\end{tabular} & \begin{tabular}[t]{@{}l@{}}Rp6.750--\\Rp246.750;\\nasi/lauk\\Rp23.000--\\Rp75.000\end{tabular} & Kompetitor reputasi paling kuat karena rating dan jumlah penilaian tertinggi, serta terdapat penanda \textit{Legendaris} dan Resto Bintang 5. \\ \hline
        Restoran Garuda - Adam Malik & \begin{tabular}[t]{@{}l@{}}Rating 4,6;\\17.068 penilaian;\\jarak 6,7 km\end{tabular} & \begin{tabular}[t]{@{}l@{}}Rp14.300--\\Rp265.851;\\nasi/kotak\\Rp19.048--\\Rp92.063\end{tabular} & Jumlah penilaian jauh lebih besar dan jaraknya hampir sama dengan Nasi Gerilya, sehingga konsistensi kualitas menjadi pembeda penting. \\ \hline
        \bottomrule
    \end{tabularx}
    \tablesource{Diolah peneliti berdasarkan observasi GrabFood Nasi Gerilya tanggal 13 Juni 2026 dan observasi kompetitor tanggal 19 dan 20 Juni 2026.}
\end{table}

\begin{table}[!htbp]
    \centering
    \caption{Format Observasi Lanjutan untuk Triangulasi Waktu}
    \label{tab:format-observasi-lanjutan-grabfood}
    \footnotesize
    \setlength{\tabcolsep}{3pt}
    \renewcommand{\arraystretch}{0.95}
    \begin{tabularx}{\textwidth}{|Z{2.2cm}|Y|Y|Y|}
        \toprule
        \textbf{Tanggal / Waktu} & \textbf{Merchant} & \textbf{Aspek yang Dicek Ulang} & \textbf{Catatan Perubahan} \\ \hline
        \midrule
        Diisi saat observasi & Nasi Gerilya & Rating, jumlah penilaian, status toko, promo, rentang harga, menu tidak tersedia, dan ulasan terbaru. & Diisi perubahan dari observasi 13 Juni 2026. \\ \hline
        Diisi saat observasi & Restoran Paripurna & Rating, jumlah penilaian, jarak, promo, harga, menu, dan ulasan terbaru. & Diisi perubahan dari observasi 19 Juni 2026. \\ \hline
        Diisi saat observasi & Rumah Makan Dendeng Batokok & Rating, jumlah penilaian, jarak, promo, harga, menu, dan ulasan terbaru. & Diisi perubahan dari observasi 19 Juni 2026. \\ \hline
        Diisi saat observasi & Rumah Makan Istana Krakatau & Rating, jumlah penilaian, jarak, promo, harga, menu premium, dan ulasan terbaru. & Diisi perubahan dari observasi 20 Juni 2026. \\ \hline
        Diisi saat observasi & RM. Pondok Gurih & Rating, jumlah penilaian, penanda reputasi, jarak, promo, harga, menu, dan ulasan terbaru. & Diisi perubahan dari observasi 20 Juni 2026. \\ \hline
        Diisi saat observasi & Restoran Garuda & Rating, jumlah penilaian, status toko, jarak, promo, harga, menu, dan ulasan terbaru. & Diisi perubahan dari observasi 20 Juni 2026. \\ \hline
        Diisi saat observasi & Kompetitor tambahan di luar lima pembanding awal & Merchant lain yang disebut pemilik atau pelanggan. & Diisi jika wawancara menunjukkan kompetitor yang lebih relevan. \\ \hline
        \bottomrule
    \end{tabularx}
    \tablesource{Diolah peneliti sebagai format observasi lanjutan (2026).}
\end{table}
